\documentclass[aspectratio=169,10pt]{beamer}
\usetheme[numbering=fraction,progressbar=frametitle]{metropolis}
\usepackage{fontspec}
\setsansfont{Latin Modern Sans}
\usepackage{amsmath,amssymb,mathtools,graphicx,booktabs}
\usepackage{polyglossia}
\setdefaultlanguage{english}
\setbeamertemplate{navigation symbols}{}
\setbeamerfont{frametitle}{size=\large}
\setbeamerfont{normal text}{size=\normalsize}
\setlength{\parskip}{.45em}
\hypersetup{pdfauthor={},pdfsubject={Lecture slides},pdfkeywords={},pdftitle={Coordinate geometry: vectors and lines}}
\providecommand{\N}{\mathbb{N}}
\providecommand{\Z}{\mathbb{Z}}
\providecommand{\D}{\mathbb{D}}
\providecommand{\Q}{\mathbb{Q}}
\providecommand{\R}{\mathbb{R}}
\providecommand{\C}{\mathbb{C}}
\providecommand{\vect}[1]{\overrightarrow{#1}}
\providecommand{\vi}{\vec{\imath}}
\providecommand{\vj}{\vec{\jmath}}
\providecommand{\vk}{\vec{k}}
\providecommand{\norm}[1]{\left\lVert #1\right\rVert}
\providecommand{\vc}[2]{\begin{pmatrix}#1\\#2\end{pmatrix}}
\providecommand{\deter}[2]{\det\left(#1,#2\right)}
\providecommand{\abs}[1]{\left\lvert #1\right\rvert}
\providecommand{\dd}{\,\mathrm{d}}
\providecommand{\e}{\mathrm{e}}
\providecommand{\Card}{\operatorname{Card}}
\providecommand{\E}{\operatorname{E}}
\providecommand{\V}{\operatorname{V}}
\providecommand{\degre}{^{\circ}}
\providecommand{\un}[1]{\,\mathrm{#1}}
\providecommand{\sci}[2]{#1\times10^{#2}}
\providecommand{\pgcd}{\operatorname{PGCD}}
\providecommand{\ppcm}{\operatorname{PPCM}}
\title{Coordinate geometry: vectors and lines}
\subtitle{Lecture and worked examples}
\date{}
\author{}
\begin{document}
\begin{frame}[plain]
\vfill
{\usebeamerfont{title}\inserttitle\par}
\medskip
{\usebeamerfont{subtitle}\insertsubtitle\par}
\medskip
{\usebeamercolor[fg]{progress bar}\rule{\textwidth}{.5pt}}
\vfill
\end{frame}
\begin{frame}{Parallelogram and intersection}
\begin{columns}[T,onlytextwidth]
\begin{column}{.57\textwidth}
In an orthonormal frame, \(A(-4,-3)\), \(B(1,-1)\) and \(C(-1,3)\).

\begin{enumerate}
\setcounter{enumi}{0}
\item Determine \(D\) such that \(ABCD\) is a parallelogram.

\item Determine its centre \(I\) in two ways.

\item With \(E(2,1)\), determine \(K\), the intersection of \((AB)\) and \((CE)\).

\item Determine \(F\) such that \(\vect{AF}=\dfrac23\vect{AE}\). Are \(D,F,B\) collinear?

\end{enumerate}
\end{column}
\begin{column}{.4\textwidth}
\centering
\includegraphics[width=\linewidth,height=4.5cm,keepaspectratio]{figures/seconde-12-1-en.pdf}
\end{column}
\end{columns}
\end{frame}
\begin{frame}{Worked answer: Parallelogram and intersection}
1. \(D=A+C-B=(-6,1)\). 2. Midpoint of \([AC]\) and of \([BD]\): \(I(-5/2,0)\). 3. \((AB):2x-5y-7=0\); \((CE):2x+3y-7=0\). Subtracting gives \(8y=0\), then \(x=7/2\), so \(K(7/2,0)\). 4. \(\vect{AE}=(6,4)\), so \(F=A+(4,8/3)=(0,-1/3)\). \(\vect{DF}=(6,-4/3)\), \(\vect{DB}=(7,-2)\); determinant \(-12+28/3=-8/3\ne0\). The three points are not collinear.
\end{frame}
\begin{frame}{Comparing a vector method and a coordinate method}
In a non-degenerate triangle \(ABC\), \(I\) is the midpoint of \([AB]\) and \(J\) the midpoint of \([AC]\).

\begin{enumerate}
\setcounter{enumi}{0}\item Prove with vectors that \((IJ)\parallel(BC)\).
\item Check with \(A(0,0)\), \(B(6,0)\), \(C(2,4)\).
\item For these coordinates, determine the equations of \((IJ)\) and \((BC)\), then the ratio \(IJ/BC\).

\end{enumerate}
\end{frame}
\begin{frame}{Worked answer: Comparing a vector method and a coordinate method}
1. \(\vect{IJ}=\vect{IA}+\vect{AJ}=\dfrac12(\vect{BA}+\vect{AC})=\dfrac12\vect{BC}\). The vectors are nonzero and collinear.

2. \(I(3,0)\), \(J(1,2)\); \(\vect{IJ}=(-2,2)\) and \(\vect{BC}=(-4,4)\).

3. \((IJ):x+y-3=0\); \((BC):x+y-6=0\). Their direction vectors are collinear. \(IJ=\sqrt8\), \(BC=\sqrt{32}\), so \(IJ/BC=1/2\). The vector proof is general; the numerical check only proves this particular case.
\end{frame}
\begin{frame}{Comparing a vector method and a coordinate method}
A figure seems to show two parallel lines. What is the difference between this observation and a proof?
\end{frame}
\begin{frame}{Worked answer: Comparing a vector method and a coordinate method}
The figure suggests a conjecture and may be imprecise. A proof uses exact data and a property, for example a zero determinant of nonzero direction vectors. A conclusion must state the objects concerned.
\end{frame}
\begin{frame}{Recognising a configuration and justifying}
In an orthonormal frame, \(A(0,0)\), \(B(4,2)\), \(C(3,4)\) and \(D(-1,2)\).

\begin{enumerate}
\setcounter{enumi}{0}\item Show that \(ABCD\) is a parallelogram.
\item Calculate the four lengths and determine whether it is a rhombus.
\item Calculate \(AC^2\), \(BD^2\) and determine whether it is a rectangle.
\item Determine its centre and an equation of each diagonal.

\end{enumerate}
\end{frame}
\begin{frame}{Worked answer: Recognising a configuration and justifying}
1. \(\vect{AB}=(4,2)=\vect{DC}\): it is a parallelogram. 2. \(AB=CD=\sqrt{20}=2\sqrt5\), \(BC=DA=\sqrt5\); adjacent sides are not equal, so it is not a rhombus. 3. \(AC^2=3^2+4^2=25\); \(BD^2=(-5)^2+0^2=25\). A parallelogram whose diagonals have the same length is a rectangle. One can also use \(AB^2+BC^2=20+5=25=AC^2\) and the converse of Pythagoras. 4. The centre is \(I(3/2,2)\). \((AC):4x-3y=0\); \((BD):y=2\). Their intersection recovers \(I\).
\end{frame}
\end{document}
